\scanpage{50}

This could be given in an extra \unsure{paragraph} (which I did not really touch upon in the talk), containing also the deduction of the Weil conjectures from the conjectures $C$ and $A$.

A last and rather trivial remark is the following. Let's introduce \deleted{a} variants $A'_{\unsure{\ell|\chi}}(X)$ and $A''_{\unsure{\ell|\chi}}(X)$ as follows:

$A'_{\unsure{\ell|\chi}}(X):$ if $2i\leq n-1$, any element \inserted{$x$} of $H^i(X)$ whose image in $H^i(Y)$ is algebraic, is algebraic.

$A''_{\unsure{\ell|\chi}}(X):$ if $2j\geq n-1$, any \inserted{algebraic} element \deleted{\unsure{of} $H^i(Y)$ whose image in} \inserted{of} $H^{j+2}(X)$ is \deleted{algebraic, is algebraic \unsure{is}} the image of an algebraic element of $H^i(Y)$.

Let \unsure{us} consider also the specifications $A'_{\unsure{\ell|\chi}}(X)^0$ and $A''_{\unsure{\ell|\chi}}(X)^0$, where we restrict to the \unsure{crittcal|critical} dimensions $2i=n-1$ if $n$ odd, $2i=n-2$ if $n$ even. All these conditions are in the nature of ``week'' Lefschets relations, and they are trivially implied by $A_{\unsure{\ell|\chi}}(X)$ resp. $C_{\unsure{\ell|\chi}}(Y)$ (in the first case, applying $\phi$ we see that $L_Xx$ is algebraic; in the second, we take $y=\Lambda_Y\phi^{\unsure{+|*}}(x)$). The remark then is that these pretend\inserted{ed}ly ``week'' variants in fact imply the full Lefschets relations for al\inserted{ge}braic cycles, namely:

\uline{Proposition}\quad $\unsure{C|A}_{\unsure{\ell|\chi}}(X)$ is equivalent to the \unsure{conjunction} $C_{\unsure{\ell|\chi}}(Y)+A'_{\unsure{\ell|\chi}}(X\times X)^0+A''_{\unsure{\ell|\chi}}(X\times X)^0$, hence (by induction) also to the conjunction of the conditions ${A'_{\unsure{\ell|\chi}}}^0$ and ${A''_{\unsure{\ell|\chi}}}^0$ for all of the varieties $X\times X,Y\times Y,Z\times Z\ \mathinner{\ldotp\ldotp\ldotp\ldotp}\ .$ Analogous statement with $X\times Y,Y\times Z$ etc instead of $X\times X,Y\times Y$ etc.

This comes from the remark that \deleted{if $X$ is even-dimensional (say a square like $X\times X$) then} $A_{\unsure{\ell|\chi}}(X\deleted{\unsure{\times X}})^0$ follows from the conjunction of $A'_{\unsure{\ell|\chi}}(X)^0$ and $A''_{\unsure{\ell|\chi}}(X)^0$, as one sees by decomposing $L_X^2:H^{2m-2}(X)\to H^{2m+2}(X)$ into
\diagram{\begin{tikzcd}[ampersand replacement=\&,column sep=small]
H^{2m-2}(X)\arrow[r,"\varphi^*"] \& H^{2m-2}(Y)\arrow[r,"\varphi_*"] \& H^{2m}(X)\arrow[r,dash,"L_X"] \& H^{2m+2}\deleted{\unsure{X}}(X)
\end{tikzcd}}
if $\dim X=2m$ is even, and $H^{\unsure{2m+1-1}}(X)\to H^{2m+1+1}(X)$ into
\diagram{\begin{tikzcd}[ampersand replacement=\&,column sep=small]
H^{2m}(X)\arrow[r,dash,"\varphi^*"] \& H^{2m}(Y)\arrow[r,dash,"\varphi_*"] \& H^{2m+2}(X)
\end{tikzcd}}
if $\dim X=2m+1$ is odd.

Sincerely your's

\mspage{49}
